% !TEX root = ../main.tex

\chapter{App Development and Testing}

\section{Frozen Revisions, Environment, and Reproducibility}

This chapter evaluates the system against inspected source revisions rather than design documents.
Two revisions bound the claims that follow.
The integrated repository state is commit \texttt{b945d3c} of the MNEME repository --- the same snapshot Chapter~3 declares as its implementation cutoff --- which contains the backend service, the AI service layer, the behavior baseline, the citation-graph platform, the integrated Android client, and their test suites.
An earlier draft of this chapter cited the older snapshot \texttt{70eb54c}; the AI service tree (\texttt{backend/src/mneme/ai} and its configuration) was verified identical between the two revisions by direct comparison, so no AI implementation claim changed under the alignment, while backend and client rows in Table~\ref{tab:status} now describe the integrated state.
The AI evaluation additions of Milestone~4 --- the graded evaluation harness, the expanded fixture set, the live evaluation command, and the cost audit --- exist on branch \texttt{feat/ai-m4-eval} (tip \texttt{59af174}, pull request open at the time of writing) and are cited with that revision explicitly.
A statement about code that appears in neither revision is a design statement and belongs to Chapter~3.

The backend is a Python/FastAPI service with PostgreSQL and \texttt{pgvector} for storage and vector search, Redis for caching, budget accounting, and job queues, and ARQ workers for background processing.
The AI layer routes each task to a configured provider; Anthropic, OpenAI, and DeepSeek adapters are implemented, providers are inferred from model identifiers, and misconfigured routes fail at startup rather than on the first request.
Embeddings default to \texttt{text-embedding-3-small} at 1536 dimensions, matching the frozen \texttt{vector(1536)} schema; each chunk row stores its embedding model so a model swap cannot leave orphan vectors.
Every completion persists provider, model snapshot, prompt version, input hash, token counts, estimated cost, and latency, so any live run is auditable from the database without extra instrumentation.
A deterministic fake provider drives the software tests; no software test in this thesis is presented as a research result.

\section{Development Process and Interim Deliverables}

\subsection{Milestone Decomposition}

Development was organized around vertical evidence paths rather than completing isolated screens.
The first path establishes paper identity, metadata ingestion, durable jobs, and a protected catalog response.
The second prepares a selected revision into sections, chunks, summaries, and embeddings.
The third exposes source-linked retrieval and Q\&A.
The fourth connects explicit interests and behavioral events to a versioned briefing.
The fifth presents a bounded graph and returns mobile events to the server.
This ordering limits late integration risk because every additional mechanism must produce a stored result, a public response, a client representation, a visible state, and an executable check.

The mid-term report records completion of foundational FastAPI, persistence, AI-service, and Android state artifacts.
It also records deliberate gaps: final connected briefing refresh, notification UX, complete graph recovery, final live quality runs, and broader usability evidence.
Those gaps are retained in the status table rather than being hidden by the presence of a prototype.

\subsection{Backend and AI Deliverables}

At the interim cutoff, the backend exposes session, catalog, job, summary, question, interest, event, briefing, and graph boundaries.
Paper records and observed revisions are separated; workers perform long-running preparation; and job state survives process restart.
The AI layer supplies deterministic test adapters, provider routing, prompt and model identity, token and cost telemetry, summary schemas, chunk embeddings, single-paper retrieval, bounded evidence markers, refusal, and limited source matching.
The data model stores the lineage needed to connect a completion back to its paper revision and source chunks.

These deliverables demonstrate software feasibility and create research instrumentation.
They do not establish parser quality, answer correctness, recommendation benefit, or user trust.
Those claims require E1--E3 and the final E4 workflow respectively.

\subsection{Android Deliverables}

The Android client is organized around Compose presentation, ViewModels, repositories, Retrofit network services, Room persistence, and WorkManager background execution.
The implemented state model distinguishes loading, live, cached, controlled fixture, empty, fallback, and retryable error paths where applicable.
Briefings and paper records can be retained locally; pending behavioral events use UUIDs and retry metadata; the graph uses a local D3 bundle inside a WebView while selection is mirrored into native state.

Heng's implementation scope covers the live briefing and paper states backed by local cache, the recoverable event-delivery path, background refresh/notification integration, and the bounded citation-graph client.
At the stated cutoff, cache-backed states, event queueing, WorkManager recovery, and graph interaction are present in code; live briefing refresh and the complete notification permission/delivery flow remain planned.
This split is repeated in the E4 matrix so that ownership does not become an unqualified completion claim.

\begin{figure}[htbp]
  \centering
  \includegraphics[page=5,width=0.94\textwidth]{figures/source_materials/uiux_mockups.pdf}
  \caption{Read--Ask--Connect branches used to align Android screens with paper, Q\&A, graph, search, and interest resources during development.}
  \label{fig:read-ask-connect}
\end{figure}

\begin{figure}[htbp]
  \centering
  \includegraphics[page=8,width=0.94\textwidth]{figures/source_materials/storymap_architecture.pdf}
  \caption{End-to-end data flow from ingestion through personalization, evidence-linked engagement, and citation-graph connection in the supplied engineering pack.}
  \label{fig:development-flow}
\end{figure}

\subsection{Testing Layers}

The project uses four evidence layers.
Unit tests check deterministic transformations and boundary validation.
Repository and API tests check transactions, identity, idempotency, and controlled responses.
Integration tests check that client and server models agree and that visible state follows backend outcomes.
Research evaluations measure quality, usefulness, latency, cost, and failure on frozen datasets or participant tasks.
Passing a lower layer is necessary but never substituted for a higher layer.

\begin{table}[htbp]
\centering
\caption{Development evidence layers and the conclusions each layer permits.}
\label{tab:test-layers}
\begin{tabular}{P{2.7cm}P{5.2cm}P{5.4cm}}
\toprule
Layer & Example artifact & Permitted conclusion \\
\midrule
Unit & parser helper, scorer, verifier, state reducer test & the function follows its specified rule on the test inputs \\
Repository/API & revision isolation, idempotent job/event, response validation & the boundary preserves stated identity and transaction invariants \\
Integration/device & client mapping, Room migration, worker retry, graph bridge & components interoperate in the recorded environment and failure case \\
Research evaluation & frozen corpus, QA labels, replay set, moderated task data & measured quality or usefulness within the stated sample and configuration \\
\bottomrule
\end{tabular}
\end{table}

\begin{figure}[htbp]
  \centering
  \includegraphics[page=14,width=0.94\textwidth]{figures/source_materials/midterm_report.pdf}
  \caption{Milestone~1 engineering-result summary from the dated mid-term report. It is retained as an interim-cutoff artifact and does not substitute for the final frozen evaluation.}
  \label{fig:midterm-engineering-results-1}
\end{figure}

\begin{figure}[htbp]
  \centering
  \includegraphics[page=15,width=0.94\textwidth]{figures/source_materials/midterm_report.pdf}
  \caption{Continuation of the Milestone~1 implementation and test inventory, including backend and Android evidence available at that cutoff.}
  \label{fig:midterm-engineering-results-2}
\end{figure}

\begin{figure}[htbp]
  \centering
  \includegraphics[page=16,width=0.94\textwidth]{figures/source_materials/midterm_report.pdf}
  \caption{Interim artifact and verification matrix used to distinguish executable software evidence from pending research measurements.}
  \label{fig:midterm-artifact-matrix}
\end{figure}

\begin{figure}[htbp]
  \centering
  \includegraphics[page=17,width=0.94\textwidth]{figures/source_materials/midterm_report.pdf}
  \caption{Reproducibility and result-traceability requirements documented at the mid-term checkpoint.}
  \label{fig:midterm-reproducibility}
\end{figure}

\section{Evidence and Implementation Status}

Table~\ref{tab:status} classifies the current evidence.
The classification prevents design intent from being promoted into an implementation claim and prevents implemented software from being promoted into a measured research outcome.

\begin{table}[htbp]
\centering
\small\renewcommand{\arraystretch}{0.92}
\caption{Current evidence status at the stated revisions. Code presence is verified by inspection; behavioral claims for backend and client subsystems require their domain owners' sign-off.}
\label{tab:status}
\begin{tabular}{P{4.5cm}P{3.2cm}P{6.3cm}}
\toprule
Artifact or claim & State & Basis and boundary \\
\midrule
AI service layer: section-aware chunking, embeddings, dense retrieval, rerank, grounded Q\&A with refusal and citation verification, summarization, recommendation scoring, caching, budget guard, model routing & \implemented{} & \texttt{b945d3c}, \texttt{backend/src/mneme/ai}; unit and API tests pass with a fake provider. Tests establish specified behavior, not answer quality. \\
Cross-revision retrieval isolation & \implemented{} (software test) & Chunks bind to \texttt{paper\_version\_id}; an API test asserts Q\&A retrieval cannot mix revisions. E1 research metrics remain \pending. \\
RAG-graded evaluation harness, 15-case fixture set, live evaluation command & \implemented{} & \texttt{feat/ai-m4-eval} at \texttt{59af174}; not yet merged into the integrated branch. \\
E2 quality results (recall, source match, refusal, latency, cost) & \pending & Requires a live run of the evaluation command against ingested fixture papers. \\
AI cost audit & \observed{} & Static call-site and list-price analysis at \texttt{1a604b3}; estimates, not measured spend. \\
Backend ingestion, jobs, behavior-v1, and graph platform & \implemented{} (code) & Present with tests at \texttt{b945d3c}; behavioral claims pending Ruiyu's sign-off. E1/E3 research results \pending. \\
Android client, offline cache, event queue, graph UI & \implemented{} (code) & Present with tests at \texttt{b945d3c}; behavioral and device claims pending Heng's and Hanyang's sign-off. \\
Adaptive recommendation effectiveness & \pending & behavior-v1 supplies the behavior embedding; E3 event streams, baselines, and results absent. \\
Mobile latency and reliability & \pending & Device matrix and system measurements absent. \\
Five-participant prototype findings & \observed{}; \limited{} & Task counts appear in the final presentation; participant-level records absent. \\
\bottomrule
\end{tabular}
\end{table}

\section{E1: Revision and Representation Evaluation}

\subsection{Evaluation Question}

E1 tests the contribution promise that every claim-bearing artifact can be traced to the intended paper revision and that a revision update cannot silently leave stale derivatives marked current.

\subsection{Implemented Basis}

Two mechanisms from the inspected revision already constrain this evaluation.
Chunking is deterministic: the same parsed sections always produce the same chunk contents and content hashes, which makes the embed stage idempotent and makes re-chunking comparable across runs.
Retrieval is revision-scoped: each chunk row binds to a \texttt{paper\_version\_id}, and an implemented API test asserts that Q\&A over one revision cannot return chunks of another.
These are software facts at \texttt{b945d3c}; they do not by themselves establish parser structure quality or invalidation completeness.

\subsection{Dataset and Fixtures}

The final evaluation should freeze a corpus containing:
\begin{itemize}
  \item papers with at least two known revisions;
  \item unchanged, locally edited, and structurally changed sections;
  \item single- and multi-column layouts;
  \item equations, tables, figures, appendices, and broken text layers;
  \item metadata-only or abstract-only degraded cases.
\end{itemize}
For each fixture, a gold manifest records logical identity, revision identity, expected sections, changed regions, and affected derived artifacts.
Fixture construction and result sign-off belong to the backend/data owner.

\subsection{Metrics}

E1 reports logical-identity accuracy, revision-identity accuracy, changed-section detection, stale-artifact invalidation precision and recall, section-boundary F1, reading-order error rate, and degraded-mode accuracy.
Performance must be stratified by layout and revision-change type.
Failure analysis should identify identity collision, missed revision, parse order, boundary, caption, and invalidation errors.

\begin{table}[htbp]
\centering
\caption{E1 result schema. Values must be inserted from frozen fixtures and logs.}
\begin{tabular}{P{4.6cm}P{2.2cm}P{2.2cm}P{3.8cm}}
\toprule
Measure & Baseline & MNEME & Evidence artifact \\
\midrule
Logical/revision identity accuracy & \pending & \pending & Fixture manifest + test log \\
Stale-artifact invalidation precision/recall & \pending & \pending & Dependency trace \\
Section boundary F1 & \pending & \pending & Gold section annotations \\
Degraded-mode accuracy & \pending & \pending & Parser status log \\
\bottomrule
\end{tabular}
\end{table}

\section{E2: Retrieval, Answer, and Citation Evaluation}

\subsection{Evaluation Questions}

E2 asks whether the implemented section-aware, evidence-bounded pipeline finds the evidence a question needs, whether generated claims remain matched to their cited sources, and whether the system refuses when evidence is absent.
It also asks what each pipeline stage contributes, and at what latency and cost.

\subsection{Evaluated Pipeline}

The evaluated configuration is the implemented pipeline at \texttt{b945d3c}, not an idealized design.
Retrieval embeds the question and performs dense approximate-nearest-neighbor search over one revision's chunks in \texttt{pgvector}, returning the top \(k=8\) by cosine similarity plus up to two early-document \emph{context anchors} that preserve overview context.
Reranking is an honest lexical blend rather than a learned cross-encoder: the final score is \(0.7\) times vector similarity plus \(0.3\) times the fraction of the question's content words present in the chunk, with stopwords removed.
Ties break on original retrieval order, so reranking is deterministic for identical inputs; the top \(n=4\) chunks become evidence, and anchors outside the top \(n\) are appended rather than displacing question-specific evidence.

Generation is refused before any model call when no evidence reaches a minimum score of \(0.25\); the model may also decline by emitting a designated insufficient-evidence marker, and both paths return the same fixed refusal answer.
After generation, citation verification checks every bracketed marker against the evidence list: markers outside the list are dropped rather than repaired, and each cited chunk is compared with the local claim sentences that carry its marker.
A citation is \emph{source-matched} when at least \(30\%\) of the claim's content words appear in the cited chunk; answers are then labeled matched, partial, or insufficient.
This is retrieved-source matching, not semantic entailment, and the thesis does not claim otherwise; the distinction follows the citation-quality framing of ALCE~\cite{gao2023alce} and the hallucination evidence in RAGTruth~\cite{niu2024ragtruth}.

Completions are cached in Redis (24-hour TTL for Q\&A, 7-day for summaries) with the prompt version in the cache key, so template changes invalidate cleanly.
A shared daily budget guard in Redis refuses provider calls once the configured daily cap (default \$5) is spent; the check-then-spend sequence can overshoot by at most the in-flight completions, which is acceptable for a daily stop.

\subsection{Frozen Question Set}

The current frozen set is \texttt{qa-seed-v2} on \texttt{feat/ai-m4-eval}: 15 manually curated cases over seven arXiv papers, comprising 12 answerable questions and 3 out-of-scope cases whose correct behavior is refusal.
Each case records the target paper, a section hint, the question, a reference answer, expected keywords, and whether refusal is expected; duplicate identifiers are rejected at load time.
The set is a seed, not a benchmark: it is large enough to exercise every pipeline path and grading rule, and too small to support significance claims.
A final evaluation should extend it with multi-section, comparison, follow-up, and distractor-heavy strata, keeping questions and evidence split by paper to reduce leakage; multi-turn strata follow the motivation of mtRAG~\cite{katsis2025mtrag}.

\subsection{Run Procedure}

The live evaluation command (\texttt{python -m mneme.cli.run\_qa\_eval}) grades the full retrieval and grounded-answer pipeline against ingested, parsed, and embedded fixture papers, using the same code paths as the production route.
Fixtures whose paper is missing or unembedded are reported as skipped, never silently dropped, and skips are recorded in the output report.
Each run must record the repository revision, model routes, prompt versions, fixture version, and the full per-case JSON report as its raw artifact.

\subsection{Metrics}

The graded report computes:
\begin{itemize}
  \item \textbf{retrieval hit rate at \(k\)}: whether any retrieved section title matches the fixture's section hint, checked by bidirectional containment after case folding; this is a section-level proxy, not span-level gold annotation;
  \item \textbf{source-match rate}: the fraction of answered cases whose every citation is source-matched; partially matched answers are reported separately and never folded into the headline rate;
  \item \textbf{citation rate}: answered cases carrying at least one verified citation;
  \item \textbf{refusal accuracy} on out-of-scope cases and \textbf{false-refusal rate} on answerable cases;
  \item \textbf{keyword coverage} of expected keywords, a lexical proxy for answer content;
  \item token counts, estimated cost, and latency per case and in aggregate.
\end{itemize}
Lexical proxies are acknowledged as such.
Before any headline quality claim, automatic grades must be calibrated against a human-reviewed subset, because reference-free evaluation is not itself ground truth~\cite{es2024ragas,saadfalcon2024ares}.

\subsection{Baselines and Ablations}

The planned comparisons isolate what each stage contributes:
\begin{enumerate}
  \item long-context generation without retrieval, motivated by position-sensitivity findings~\cite{liu2024lost};
  \item fixed-size chunk retrieval versus the implemented section-aware chunking;
  \item dense retrieval only versus the vector-plus-lexical rerank;
  \item the pipeline without context anchors;
  \item the pipeline without citation verification;
  \item varied refusal thresholds around the deployed \(0.25\).
\end{enumerate}
The comparisons guard against assuming that every additional stage improves quality or latency~\cite{wang2024bestpractices}.
Naming a cross-encoder or another learned reranker as part of MNEME is gated on implementing it and passing this ablation; until then, vector-plus-lexical remains the accurate description.

\begin{table}[htbp]
\centering
\caption{E2 result schema. Component metrics prevent a fluent answer from hiding retrieval or citation failure.}
\begin{tabular}{P{4.0cm}P{2.1cm}P{2.1cm}P{2.1cm}P{2.1cm}}
\toprule
Configuration & Retrieval & Source match / refusal & Keyword coverage & Latency / cost \\
\midrule
Long context & \pending & \pending & \pending & \pending \\
Fixed chunks & \pending & \pending & \pending & \pending \\
Dense only & \pending & \pending & \pending & \pending \\
Full pipeline (deployed) & \pending & \pending & \pending & \pending \\
\bottomrule
\end{tabular}
\end{table}

\section{E3: Interest-Memory Evaluation}

\subsection{Evaluation Questions}

E3 tests whether adaptive interest state improves future paper ranking under a temporal split, whether short- and long-term state serve distinct roles, and whether users can correct the model consistently.
The evaluation does not treat historical clicks as perfect preference labels.

\subsection{Implemented Scoring Basis}

The implemented recommendation scorer at \texttt{b945d3c} blends three signals: explicit topic match (weight 0.45), cosine similarity between the user's learned behavior embedding and the paper's mean chunk embedding (weight 0.35), and publication recency with a one-week half-life (weight 0.2).
Missing signals redistribute their weight over the available ones, so a cold-start user still receives a meaningful score; ranking is deterministic with a stable tie-break.
Each entry carries human-readable reasons derived from its dominant signals, not free-form post-hoc text, consistent with the scrutability requirement of Chapter~3~\cite{ramos2024scrutable}.
The behavior embedding itself is produced by behavior-v1, owned and signed off by the backend/data owner; this chapter evaluates the composition, not the internals of the behavior model.

\subsection{Data and Baselines}

The final evaluation freezes event sequences with exposure records, event semantics, timestamps, and paper candidates.
Where real longitudinal data are unavailable, synthetic sequences may test deterministic update invariants but cannot establish user benefit.
Baselines include explicit interests only, recency only, behavior-v1 with the deployed weights, and any frozen v2 candidate; a v2 comparison is reportable only after its definition and migration are frozen.

\subsection{Metrics and Tests}

Ranking metrics include Recall@\(k\) and nDCG@\(k\) on future eligible positive events.
Diversity and novelty are reported to detect narrowing or repetition.
Adaptation delay measures how quickly a confirmed new interest affects ranking.
Negative-feedback sensitivity tests whether a bounded skip signal has an appropriate rather than overwhelming effect, respecting the exposure-eligibility rule of Chapter~3~\cite{zhang2025unpaired}.
Correction consistency tests whether editing or deleting a visible interest produces predictable ranking and explanation changes.

\begin{table}[htbp]
\centering
\caption{E3 result schema. Metrics must be reported on a temporal split with exposure-aware labels.}
\begin{tabular}{P{4.0cm}P{2.1cm}P{2.1cm}P{2.1cm}P{2.1cm}}
\toprule
Model & Ranking & Diversity & Adaptation & Correction \\
\midrule
Explicit only & \pending & \pending & \pending & \pending \\
Recency only & \pending & \pending & \pending & \pending \\
behavior-v1 (deployed weights) & \pending & \pending & \pending & \pending \\
Frozen v2 candidate (if any) & \pending & \pending & \pending & \pending \\
\bottomrule
\end{tabular}
\end{table}

\section{E4: Systems and Usability Evaluation}

\subsection{Systems Measurements}

The systems evaluation requires an immutable build, backend revision, model and provider configuration, device matrix, network conditions, and corpus snapshot.
It reports digest-load latency, paper-detail latency, Q\&A first-token and completion latency, graph expansion latency, event-upload success, synchronization conflicts, background-task completion, notification delivery, cache hit rate, and daily model cost.
Each distribution includes sample count and percentile statistics.
The AI-side instrumentation for this evaluation is already implemented: every completion persists token counts, estimated cost, and latency, and budget exhaustion surfaces to the client as an explicit rate-limit response rather than a silent failure.
Failure injection covers offline operation, provider timeout, budget exhaustion, malformed document, empty retrieval, verification failure, and duplicate event upload.
Client and device measurements are supplied and signed off by the client owners.

\subsection{Usability Protocol}

The supplied task design is retained because it covers the core research loop:
\begin{enumerate}
  \item identify why a paper was recommended and open it;
  \item identify the main contribution and verify a source;
  \item ask a dataset question and inspect cited evidence;
  \item open the graph and inspect a related paper;
  \item change a digest preference or interest topic.
\end{enumerate}
The final dataset must contain one row per participant and task with completion, time, taps, error path, confusion code, and observation.
The study report must describe recruitment, consent, moderator instructions, prototype or build version, analysis procedure, and missing data.

The moderator introduction states that the interface, not the participant, is being tested and asks the participant to think aloud.
The facilitator does not teach where to click.
The prototype protocol defines expected performance targets: T1 within 20 seconds and at most two taps; T2 within 45 seconds with a source opened; T3 within 60 seconds with an answer and citation found; T4 within 45 seconds and at most three taps; and T5 within 60 seconds with a saved preference.
These targets are design criteria.
Because the supplied aggregate report does not include per-participant time and tap records, the present thesis reports completion counts only.

\begin{figure}[htbp]
  \centering
  \includegraphics[page=10,width=0.94\textwidth]{figures/source_materials/uiux_mockups.pdf}
  \caption{Prototype tasks, expected actions, and success signals. Timing and tap targets remain criteria until participant-level records are supplied.}
  \label{fig:usability-protocol}
\end{figure}

\subsection{Available Formative Observation}

As stated in Chapter 2, the final presentation reports task-success counts of 5, 5, 5, 2, and 4 out of five.
Within that report, the graph task is the only major drop-off.
This observation supports the decision to make graph actions explicit.
It does not show that the subsequent redesign improves completion because no retest is supplied.

\subsection{Android Evaluation Scope}
\label{sec:e4-android}

The executed Android subset of E4 measured one Android 14/API 34 Google emulator with four virtual processors and a 192~MiB application heap.
Three startup pairs were discarded as warm-up.
The retained startup set contains 20 cold process starts and 20 task foreground/resume operations.
Controlled state, graph, and event tests used their own discarded warm-up iterations.
Medians are sample medians and p95 values use the nearest-rank definition.
Appendix~\ref{app:android-e4} gives the complete sample inventory and success conditions.

The repeated trials cover process and task activation, five content-state conditions, one polling path, four graph sizes, three graph disclosures, and four controlled event stages.
All samples met the success condition defined for their group; their audit total is 810.
The groups remain separate in the results below.
The live-backend trace contains five repetitions with five recorded stages per repetition; it is not pooled with the controlled results.
Together, the subsections map the acceptance conditions in Table~\ref{tab:mobile-acceptance} to observed client behavior.
State and polling results address content and asynchronous-work conditions, graph results address selected-state and navigation conditions, and event results address retry and duplicate acknowledgement.
Because no participant retest was run, this mapping establishes software acceptance only.

\subsection{Mobile State and Reliability Results}

Figure~\ref{fig:mobile-reliability} and Table~\ref{tab:mobile-reliability} report startup, resume, controlled state, and polling latency.
All 20 cold starts and all 20 task foreground/resume operations returned a wait-time record for the requested Android activity.
The startup runner does not inspect rendered UI content.
Cold process start had a median of 2,773~ms and p95 of 4,444~ms.
The task foreground/resume operation had a median of 133.5~ms and p95 of 445~ms.
Android classified nine operations as hot starts and returned \texttt{UNKNOWN (0)} for eleven.
The distribution uses the retained wait-time field and is limited to the foreground/resume operation.
The records do not constitute 20 independently confirmed hot starts.

\begin{figure}[htbp]
\centering
\includegraphics[width=0.92\textwidth]{mneme_mobile_reliability}
\caption{Android client latency and reliability. Separate panels prevent process/task timing, controlled content-state transitions, and event operations from being treated as one timing family. Controlled state timings begin after fixture input; they are not network or display-frame latency.}
\label{fig:mobile-reliability}
\end{figure}

\begin{table}[htbp]
\centering
\caption{Android startup, controlled state, and polling results. Latencies are in milliseconds.}
\label{tab:mobile-reliability}
\begin{tabular}{P{5.1cm}P{1.5cm}P{2.2cm}P{2.2cm}P{2.2cm}}
\toprule
Condition & \(n\) & Successful & Median & p95 \\
\midrule
Cold process start & 20 & 20 & 2,773 & 4,444 \\
Task foreground/resume & 20 & 20 & 133.5 & 445 \\
\addlinespace
Live seed transition (controlled) & 30 & 30 & 0.124 & 0.293 \\
Cached warm start (controlled) & 30 & 30 & 0.227 & 0.376 \\
Cached offline disclosure (controlled) & 30 & 30 & 0.227 & 0.439 \\
Cached server-failure disclosure (controlled) & 30 & 30 & 0.231 & 0.362 \\
Invalid-cache error (controlled) & 30 & 30 & 0.203 & 0.635 \\
\addlinespace
Three-stage job polling (controlled) & 10 & 10 & 3,008.418 & 3,064.345 \\
\bottomrule
\end{tabular}
\end{table}

The five controlled content-state conditions each completed 30 of 30 trials.
Their sub-millisecond values measure the ViewModel transition after a deterministic fixture result is delivered.
The values cannot be interpreted as live service latency or UI frame time.
The three-stage paper job completed in all ten trials with a median of 3,008.418~ms, consistent with the one-second client polling interval across three processing stages.

\subsection{Bounded Graph Results}

The graph experiment exercised 1, 12, 25, and 50 nodes.
The upper bound matches the client request contract.
Each size was measured in a newly created WebView and as an update to an existing WebView.
Figure~\ref{fig:graph-latency} shows the resulting distributions; Table~\ref{tab:graph-latency} reports their medians and p95 values.

\begin{figure}[htbp]
\centering
\includegraphics[width=0.92\textwidth]{mneme_graph_latency}
\caption{Citation-graph presentation and selection latency by node count. Render timing ends when the local graph page reports DOM readiness. It does not measure convergence of the force-directed layout.}
\label{fig:graph-latency}
\end{figure}

\begin{table}[htbp]
\centering
\caption{Graph DOM-ready and native selection-callback latency. Each cell is median / p95 in milliseconds.}
\label{tab:graph-latency}
\begin{tabular}{P{1.3cm}P{3.0cm}P{3.0cm}P{3.0cm}P{3.0cm}}
\toprule
Nodes & New WebView ready & Existing WebView ready & New WebView selection & Existing WebView selection \\
\midrule
1  & 139.749 / 298.019 & 83.001 / 150.840 & 32.764 / 46.022  & 32.279 / 103.680 \\
12 & 454.149 / 913.746 & 95.590 / 331.895 & 35.971 / 392.792 & 32.968 / 208.816 \\
25 & 315.039 / 734.190 & 97.767 / 361.523 & 48.675 / 568.963 & 32.759 / 128.642 \\
50 & 306.793 / 928.176 & 111.814 / 513.785 & 49.098 / 353.710 & 37.860 / 163.644 \\
\bottomrule
\end{tabular}
\end{table}

Every ready and selection sample succeeded at each evaluated size.
Median DOM-ready latency ranged from 139.749 to 454.149~ms for a new WebView and from 83.001 to 111.814~ms for an existing WebView.
Median selection-callback latency ranged from 32.279 to 49.098~ms.
The non-monotonic medians and wider p95 values at several sizes show emulator scheduling and WebView initialization variability; the experiment does not support a linear scaling model.

Ready, fallback, and empty graph states each completed 30 of 30 controlled disclosure trials.
The existing graph and app-flow suite passed 11 of 11 tests, including node selection, Open Paper navigation, and selection restoration after returning from paper detail.
These results establish the specified presentation and navigation behavior.
They do not measure graph usefulness, discoverability, ranking, clustering, or layout quality.

\subsection{Event Synchronization and Live Readback}

The controlled event experiment isolated local queue and acknowledgement behavior from network variability.
Table~\ref{tab:event-controlled} reports 30 trials for each stage.
All queue writes, retryable failures, subsequent retries, and duplicate replays met their success conditions.

\begin{table}[htbp]
\centering
\caption{Controlled client event-synchronization results. Latencies are in milliseconds.}
\label{tab:event-controlled}
\begin{tabular}{P{5.0cm}P{1.3cm}P{2.0cm}P{2.0cm}P{2.0cm}}
\toprule
Stage & \(n\) & Successful & Median & p95 \\
\midrule
Durable queue write & 30 & 30 & 10.287 & 32.261 \\
Retryable HTTP failure & 30 & 30 & 15.613 & 80.332 \\
Retry success & 30 & 30 & 54.999 & 91.794 \\
Duplicate replay & 30 & 30 & 18.950 & 50.222 \\
\bottomrule
\end{tabular}
\end{table}

The live trace then exercised the same client coordinator against a running backend.
In each of five repetitions, the client queued three new events, recovered from an injected retryable failure, uploaded the batch, replayed the same identifiers, and requested a briefing through the standard client data path.
Table~\ref{tab:event-live} keeps these network-dependent stages separate from the controlled results.

\begin{table}[htbp]
\centering
\caption{Live-backend event-loop results over five repetitions. Latencies are in milliseconds.}
\label{tab:event-live}
\begin{tabular}{P{5.0cm}P{1.3cm}P{2.0cm}P{2.0cm}P{2.0cm}}
\toprule
Stage & \(n\) & Successful & Median & p95 \\
\midrule
Queue three client events & 5 & 5 & 13.850 & 19.061 \\
Injected failure before retry & 5 & 5 & 8.028 & 8.616 \\
Upload after retry & 5 & 5 & 33.049 & 188.478 \\
Replay the same identifiers & 5 & 5 & 24.210 & 125.869 \\
Read back a live briefing & 5 & 5 & 26.776 & 135.906 \\
\bottomrule
\end{tabular}
\end{table}

Each live upload accepted three new identifiers, and each replay recognized the same three identifiers as duplicates.
Every readback carried a live-backend origin.
The returned digest contained zero papers in all five repetitions, and the preference-model version remained 1.
The local evaluation database had no paper inside the active 14-day digest window.
The live result therefore supports client-loop reachability and duplicate-aware ingestion.
It supplies no evidence of preference change, recommendation adaptation, or recommendation quality.

\subsection{E4 Claim Boundary}

The device measurements come from one emulator and do not estimate physical-device variability.
Graph render timing ends at a DOM-ready callback, while the force simulation may continue.
Controlled state and event success conditions establish the specified software transitions; no participant data show whether readers understand the labels or recovery actions.
The live trace uses one local backend configuration and returns an empty digest.
Within these limits, E4 provides measured evidence for the bounded Android interaction path.
\begin{table}[htbp]
\centering
\caption{Reported formative prototype outcomes and bounded interpretations.}
\label{tab:formative-usability}
\begin{tabular}{P{0.8cm}P{5.0cm}P{1.5cm}P{5.7cm}}
\toprule
Task & Expected outcome & Reported & Bounded interpretation \\
\midrule
T1 & open the top recommended paper from the digest & 5/5 & digest-to-paper navigation was discoverable in this prototype sample \\
T2 & identify the main contribution and inspect one source & 5/5 & source-linked summary path was usable in the reported sessions \\
T3 & ask a paper question and inspect cited evidence & 5/5 & cited Q\&A supported the assigned task; answer correctness was not independently evaluated \\
T4 & select and open a related/citing paper in the graph & 2/5 & the graph lacked a clear next action and required redesign \\
T5 & change a digest preference or interest topic & 4/5 & the control worked for most participants, but its entry point required greater visibility \\
\bottomrule
\end{tabular}
\end{table}

\begin{figure}[htbp]
  \centering
  \includegraphics[page=11,width=0.94\textwidth]{figures/source_materials/thesis_uiux.pdf}
  \caption{Qualitative and quantitative synthesis reported in the UI/UX study. The thesis retains the task distribution and avoids treating the aggregate percentage as population-level evidence.}
  \label{fig:uiux-result-synthesis}
\end{figure}

\subsection{Design Changes and Retest Criteria}

The digest remains the home surface, with match information, a reason, and one dominant open action.
Verification labels and citation chips move closer to the claim they qualify.
The graph adds a concise tap instruction, an unambiguous selected-node style, a relation explanation, and a persistent ``Open paper'' button.
Interest tuning remains reachable from the digest and gains a more explicit settings entry.

\begin{figure}[htbp]
  \centering
  \includegraphics[page=8,width=0.94\textwidth]{figures/source_materials/usability_tests.pdf}
  \caption{Evidence-to-design trace retained from the usability pack: preserve the digest, strengthen verification cues, turn graph nodes into explicit actions, and keep preference tuning near the briefing.}
  \label{fig:design-change-trace}
\end{figure}

\begin{figure}[htbp]
  \centering
  \includegraphics[page=12,width=0.94\textwidth]{figures/source_materials/thesis_uiux.pdf}
  \caption{Detailed rationale for strengthening claim-adjacent verification and converting graph nodes into explicit, interpretable actions.}
  \label{fig:uiux-final-decisions}
\end{figure}

The retest should hold the assigned graph task constant and compare the revised connected build with the recorded prototype baseline.
Success requires not only task completion but also correct identification of the selected paper and relation.
The study should record false starts, taps, time, and whether the participant can explain the graph's incompleteness and fallback state.
If completion rises but participants misunderstand which paper will open, the redesign has improved discoverability without establishing semantic clarity.

\begin{figure}[htbp]
  \centering
  \includegraphics[page=10,width=0.94\textwidth]{figures/source_materials/usability_tests.pdf}
  \caption{Graph-specific failure analysis and the required instruction, selected state, relation explanation, and primary action for the next design iteration.}
  \label{fig:graph-redesign}
\end{figure}

\subsection{Mobile Reliability Matrix}

The final device evaluation should exercise each visible state under controlled conditions.
It includes first launch with no cache, live load, cache restoration, stale cache, offline operation, backend timeout, provider-budget exhaustion, empty retrieval, graph fallback, duplicate event upload, process death during upload, application restart, notification permission denial, and repeated briefing refresh.
For each case the raw artifact should contain device/OS, build revision, network condition, timestamp, expected state, observed state, event/job identity, screenshot or screen recording, and pass/fail explanation.

\begin{longtable}{P{3.2cm}P{4.2cm}P{4.1cm}P{2.0cm}}
\caption{E4 mobile reliability result schema.}
\label{tab:e4-reliability}\\
\toprule
Scenario & Expected user-visible behavior & Required raw evidence & Result \\
\midrule
\endfirsthead
\toprule
Scenario & Expected user-visible behavior & Required raw evidence & Result \\
\midrule
\endhead
Offline with cache & cached label and freshness remain visible; no new Q\&A implied & screenshot, Room snapshot, network trace & \pending \\
Offline without cache & actionable connectivity error; no fixture substitution & screenshot and state log & \pending \\
Duplicate event retry & one durable event effect and duplicate acknowledgement & queue/database trace & \pending \\
Worker interruption & reserved events return to pending or recover after timeout & WorkManager and queue log & \pending \\
Graph enrichment failure & deterministic bounded citation fallback with status label & API response and device capture & \pending \\
Notification denied & briefing remains accessible without repeated permission loop & permission and notification log & \pending \\
Repeated refresh & one notification per new briefing identity & briefing IDs and notification history & \pending \\
\bottomrule
\end{longtable}

\section{AI Cost and Latency Audit}

\subsection{Method and Scope}

Model cost is a first-class evaluation dimension because the assistant runs continuously against a paid provider.
A static audit at revision \texttt{1a604b3} inventoried every provider call site, its trigger, default model route, cache behavior, and budget check.
All dollar figures in this section are estimates from provider list prices and token approximations, classified \observed{}; measured spend comes from the persisted cost telemetry once the deployed pipeline has run, and remains \pending.
The price rows behind these figures were re-verified against the official provider pricing pages on 2026-07-26 and retained, with source links and the deployed cap assumptions, in the repository artifact \texttt{docs/evaluation/ai-cost-price-snapshot-2026-07-26.md}; the companion script \texttt{ai\_cost\_recompute.py} recomputes every dollar figure in this section from that snapshot and fails if prose and calculation diverge.

There are four provider call sites: the summarization worker (background job per paper version), grounded Q\&A (per request), query embedding, and batched chunk embedding; recommendation scoring makes no provider calls because it reuses stored vectors.
All completion sites check the shared daily budget and persist full cost telemetry.

\subsection{Estimated Unit Economics}

Under the deployed caps (summary input 60{,}000 characters, summary output 1024 tokens, Q\&A context of roughly six 450-token chunks, Q\&A output 512 tokens), one paper summary on the default flagship route costs an estimated \$0.10, one grounded answer roughly \$0.028, and embedding a full paper roughly \$0.0003.
Against the default \$5 daily cap, a single bulk ingestion day of about 50 summaries can exhaust the budget on the flagship tier, while a cheaper tier fits roughly 250 summaries per day.
A full 15-case evaluation run is estimated at \$0.34 on the flagship tier versus \$0.07 on a light tier.

\subsection{Findings}

The audit's headline finding is that the default routing sends bulk summarization --- schema-constrained background work, and the dominant spend at roughly five times Q\&A cost per unit --- to the most expensive configured model, while the flagship tier is better reserved for user-visible grounded answers.
Rerouting summarization to a cheaper tier is recommended but deliberately deferred: the decision is gated on a summary-quality comparison over the summarization fixture set, because structured-output fidelity, not only cost, decides the route.
Secondary findings: the single shared daily cap can let a heavy ingestion day starve interactive Q\&A (acceptable for the demo deployment; split per-task caps if observed); the embedding batch size is three orders of magnitude below LLM cost and needs no change; and the caching posture is sound because prompt versions key every cache entry and unknown models price at the most expensive tier, so budget accounting never undercounts.
No default was changed by the audit; the routing decision remains open with its owner recorded.

\section{Threats to Validity}

\subsection{Construct Validity}

Source matching is lexical overlap with cited evidence, not semantic entailment and not scientific truth; a span can support a claim in one paper while the broader literature disagrees~\cite{wadden2020scifact}.
Keyword coverage is a lexical proxy for answer content, and the section-hint retrieval grade is a section-level proxy for span-level relevance.
Likewise, a click or long reading duration is not identical to research interest.
The evaluation therefore reports operational constructs and avoids treating them as complete measures of understanding or intent.

\subsection{Internal Validity}

Generator, embedding, and judge models can change over time, and unrecorded provider updates may alter results.
The implemented telemetry mitigates but does not remove this risk: every persisted completion records provider, model snapshot, prompt version, and input hash, and evaluation runs must additionally freeze fixture versions and configuration.
Estimated costs derive from the retained list-price snapshot and token approximations and may diverge from billed spend.
Any LLM-as-judge measurement requires calibration against human annotations before it supports a claim.

\subsection{External Validity}

The 15-case seed set covers seven machine-learning papers from arXiv; conclusions from it do not transfer to other disciplines, layouts, or revision practices without a larger stratified corpus.
A five-participant formative prototype study cannot establish broad usability.
Results must be stratified and conclusions restricted to the evaluated documents, questions, devices, and participants.

\subsection{Reproducibility}

The source repository, revisions, fixture sets, and evaluation tooling are supplied and inspectable, which closes the largest gap of the earlier draft.
The remaining gaps are execution artifacts: no live graded run, no measured cost or latency distributions, no E1/E3 fixture corpora, and no raw usability records exist yet.
Each such gap is marked \pending{} with an owner, and the thesis becomes empirically complete only when those artifacts are produced, linked, and audited at a frozen revision.
